\section{Results}

After executing QEMU loaded with the test firmware using the command shown in 
Figure~\ref{code:qemu_test_cmd}:

\begin{figure}[H]
  \begin{lstlisting}[language=bash, caption={QEMU command for automated LIS3DH testing}, label={code:qemu_test_cmd}]
sudo ./qemu-system-arm -M netduinoplus2 -kernel firmware_test_4.elf -serial stdio
  \end{lstlisting}
\end{figure}



the following output confirms that the system fully meets all five testing objectives 
established at the beginning of the chapter:

\begin{lstlisting}[language=bash, caption={Automated test execution output}, 
label={code:test_output}]
=== LIS3DH EMULATION TEST SUITE ===

--- Test Case 1: WHO_AM_I Verification ---
 PASS: LIS3DH init OK, WHO_AM_I = 0x33

--- Test Case 2: Operating Mode Configuration ---

 Test Case 2.1: LIS3DH_MODE_NORMAL
	 TEST PASSED
 Test Case 2.2: LIS3DH_MODE_HIGH_RES
[LIS3DH] Mode: high
[LIS3DH] acc sensitivity: 1.0
	 TEST PASSED
 Test Case 2.3: LIS3DH_MODE_LOW_POWER
[LIS3DH] Mode: low
[LIS3DH] acc sensitivity: 16.0
[LIS3DH] Mode: low
[LIS3DH] acc sensitivity: 16.0
	 TEST PASSED

 PASSED:  Operating Mode Configuration Test case

--- Test Case 3: Full-Scale Configuration ---
 Test Case 3.1: LIS3DH_FS_2G
PASS
 Test Case 3.2: LIS3DH_FS_4G
[LIS3DH] FS: LIS3DH_FS_4G
[LIS3DH] acc sensitivity: 32.0
PASS
 Test Case 3.3: LIS3DH_FS_8G
[LIS3DH] FS: LIS3DH_FS_8G
[LIS3DH] acc sensitivity: 64.0
PASS
 Test Case 3.4: LIS3DH_FS_16G
[LIS3DH] FS: LIS3DH_FS_16G
[LIS3DH] acc sensitivity: 192.0
PASS

 PASS: Full-scale Test case

--- Test Case 4 & 5: Full-Scale Configuration ---
[LIS3DH] Mode: normal
[LIS3DH] acc sensitivity: 48.0
[LIS3DH] FS: LIS3DH_FS_4G
[LIS3DH] acc sensitivity: 8.0
[LIS3DH] ODR: LIS3DH_ODR_100

STM32: LIS3DH initialized with DRDY interrupt on INT1
LIS3DH: Input: 1000.0 mg -> 125.0 LSB -> 125 raw

LIS3DH: out_x_h: 0x1f, out_x_l: 0x40
LIS3DH: Input: 1000.0 mg -> 125.0 LSB -> 125 raw

LIS3DH: out_y_h: 0x1f, out_y_l: 0x40
LIS3DH: Input: 1000.0 mg -> 125.0 LSB -> 125 raw

LIS3DH: out_z_h: 0x1f, out_z_l: 0x40
LIS3DH: INT1 pulsed (DRDY, non-latched)
Calculated So = 8.000 mg/LSB (mode=0, fscale=0)
INT! X: 1000.00 mg, Y: 1000.00 mg, Z: 1000.00 mg
INT1_SRC: 0x00
\end{lstlisting}

\section{Objective-Result Mapping}

\begin{table}[H]
    \centering
    \small
    \caption{Test Objective to Result Mapping}
    \label{tab:test_objective_mapping}
    \begin{tabularx}{\textwidth}{|X|c|c|X|}
        \hline
        \textbf{Objective} & \textbf{Test Case} & \textbf{Result} & \textbf{Justification} \\
        \hline
        Protocol Fidelity & Test Case 1 & \textcolor{green}{PASSED} & 
        I$^2$C subaddressing;  auto-increment correct \\
        \hline
        Configuration Correctness & Test Case 2 \& 3 & \textcolor{green}{PASSED} & 
        Write in the control registers change the internal logic \\
        \hline
        Data Acquisition & Test Case 4 & \textcolor{green}{PASSED} & 
        QOM → registers → perfect mg scaling \\
        \hline
        Interrupt Behavior & Test Case 5 & \textcolor{green}{PASSED} & 
        DRDY → EXTI → clear non-latched working \\
        \hline
        HAL Compatibility & All (Native) & \textcolor{green}{PASSED} & 
        \texttt{HAL\_I2C\_MemRead} \texttt{/Write()} unmodified \\
        \hline
    \end{tabularx}
\end{table}

These tests, though simple in description, together with additional undocumented 
validations, verify the complete functionality of the system developed in this thesis:

\begin{itemize}
  \item QEMU LIS3DH emulator with QOM injection and 
  DRDY interrupt.
  \item STM32F405 I$^2$C controller with datasheet 
  register fidelity and native HAL support.
  \item Interrupt-driven firmware driver, portable 
  and reusable.
  \item End-to-end validation from QOM properties 
  to scaled mg output.
\end{itemize}

The results obtained demonstrate that the emulated system 
behaves identically to its physical counterpart, meeting all established objectives 
and fully validating the technical correctness of the work performed.
